% ============================================================================
% FIGURE 8: Controlled Repair Test Results
% Location: In new Section 6.4 "Controlled Repair Activation Test"
%           OR in Section 7 "Verifier Failure Analysis"
%           OR in Section 8 "Discussion"
% ============================================================================

\begin{figure}[htbp]
\centering
\includegraphics[width=0.95\textwidth]{figure8_repair_test.pdf}
\caption{Results of controlled repair activation test (N=20) where false positive
cases were forced through REWRITE by manually marking all claims as unsupported.
While repair successfully removed explicit hedging keywords ("may," "might,"
"suggests"), it replaced them with epistemic disclaimers ("remains less studied,"
"not detailed in the evidence") that still avoid definitive yes/no commitment.
Accuracy improved by only 5\% (0\% $\to$ 5\%), with 95\% of cases maintaining
"maybe" answers. This demonstrates that repair shares verification's fundamental
limitation: semantic similarity cannot detect epistemic evasion strategies that
models employ to avoid commitment.}
\label{fig:repair_test}
\end{figure}


% ============================================================================
% HOW TO ADD THIS TO YOUR REPORT
% ============================================================================

% OPTION 1: New subsection in Results (RECOMMENDED)
% Add after Section 6.3 (Ablation Studies):

\subsection{Controlled Repair Activation Test}

To isolate whether repair effectiveness is limited by verification failure or by
inherent limitations of the repair mechanism itself, we conducted a controlled
experiment. We selected 20 false positive cases where verification incorrectly
marked hedged claims as supported, manually forced all claims to be flagged as
unsupported, and measured whether REWRITE could fix the errors when properly
triggered.

Figure~\ref{fig:repair_test} shows the results. Despite successfully removing
explicit hedging keywords (100\% $\to$ 0\%), accuracy improved by only 5\%
(0\% $\to$ 5\%), with 19 out of 20 cases remaining incorrect. Analysis reveals
that while hedging keywords disappeared, they were replaced by epistemic
disclaimers (0\% $\to$ 95\%) such as ``remains less studied compared to animals''
or ``specific results are not detailed in the evidence.'' These new formulations
pass NLI verification because they maintain semantic similarity to the passages,
but they still avoid committing to definitive yes/no answers---95\% of rewritten
cases retained ``maybe'' as the answer label.

This reveals a \textbf{second failure point} beyond verification: the repair
mechanism itself cannot enforce answer commitment because it relies on the same
semantic similarity principle. Both verification and repair are vulnerable to the
epistemic evasion strategies that GPT-4o employs, demonstrating that fixing
verification alone is insufficient. The entire similarity-based approach requires
fundamental redesign to enforce definitiveness in generated claims.


% ============================================================================
% OPTION 2: Add to Discussion Section
% ============================================================================

Our controlled repair test (\S6.4, Figure~\ref{fig:repair_test}) confirms that
the problem extends beyond verification to the repair mechanism itself. When we
manually forced repair activation on 20 false positive cases, accuracy improved
by only 5\% (from 0\% to 5\%), demonstrating that REWRITE cannot fix the
underlying issue even when properly triggered. The repair mechanism successfully
removed hedging keywords but replaced them with epistemic disclaimers that still
pass NLI checks while avoiding commitment---revealing that semantic similarity is
fundamentally inadequate for enforcing answer definitiveness at \textit{both}
the verification and repair stages.


% ============================================================================
% OPTION 3: Add to Future Work
% ============================================================================

Future systems need improvements at \textit{both} verification and repair stages.
Our controlled repair test (Figure~\ref{fig:repair_test}) showed that forced
repair activation improved accuracy by only 5\% (0\% $\to$ 5\%) because the
mechanism replaced hedging keywords with epistemic disclaimers that still avoid
commitment. This suggests the need for: (1) question-aware verification that
detects epistemic evasion beyond keyword matching, and (2) repair mechanisms with
explicit anti-hedging constraints that enforce definitive yes/no answers through
iterative rewriting with verification of repaired claims.


% ============================================================================
% REFERENCES IN TEXT
% ============================================================================

% Use these patterns to reference the figure:

% "As shown in Figure~\ref{fig:repair_test}, repair barely improves accuracy..."

% "The controlled repair test (Figure~\ref{fig:repair_test}) reveals..."

% "Figure~\ref{fig:repair_test} demonstrates that repair shares verification's
%  fundamental limitation..."
